\section{Three connected research questions and verified gap}
Blockchain fee mechanisms allocate scarce inclusion capacity~\cite{roughgarden2024}, while Algorand supplies the protocol setting~\cite{gilad2017}. We study its smallest strategic fee-choice abstraction: two privately informed users simultaneously choose \emph{High} ($H$) or \emph{Low} ($L$). Equal choices create congestion and low rewards; different choices avoid the conflict. This is a one-shot static Bayesian anti-coordination game connecting equilibrium analysis, information design, and behavioral measurement.

\textbf{Q1---Economics:} How do private urgency and value affect coordination and welfare? \textbf{Q2---Computation:} Which pure, mixed, and Bayesian equilibria follow from explicit payoffs and type beliefs? \textbf{Q3---Behavior:} Does incomplete information increase same-fee failure, and can a signal reduce it? Figure~\ref{fig:teaser} connects types, choices, matrix outcomes, treatments, and metrics.

\section{Answer to Q1: incentives and strategic model}
Users $i\in\{1,2\}$ privately observe type $\theta_i\in\{U,P\}$ (urgent or patient) and simultaneously choose $a_i\in\{H,L\}$. Conditional on types,
\[
\begin{array}{c|cc}&H&L\\\hline H&(R_C,R_C)&(R_H,R_L)\\L&(R_L,R_H)&(R_C,R_C)\end{array}
\]
where $R_C$ is the same-fee congestion reward. If $R_H>R_L>R_C$, $(H,L)$ and $(L,H)$ are pure Nash equilibria~\cite{nash1950}; the symmetric mixed equilibrium chooses $H$ with $p^*=(R_H-R_C)/(R_H+R_L-2R_C)$. With private types, payoffs become type-dependent and the solution is Bayesian Nash equilibrium~\cite{harsanyi1968}. Solving it requires a type prior and type-specific payoffs, which remain calibration inputs.

\textbf{Game theory} characterizes best responses; \textbf{social choice} compares equilibrium with the welfare-maximizing assignment of one user to each fee; \textbf{mechanism design} asks whether a feasible recommendation can improve coordination without revealing types, drawing on correlated equilibrium and information design~\cite{aumann1974,bergemann2016}.

\section{Answer to Q2: computation, auction comparison, and artifacts}
The computational implementation and planned human experiment follow induced-value principles~\cite{smith1976} and compare: (1) complete information; (2) incomplete information; and (3) incomplete information plus a coordination signal. Participants would be randomly rematched across one-shot rounds. Outcomes are same-fee congestion, differentiation, urgent-user selection of $H$, welfare, and deviation from Bayesian equilibrium. The design builds on experimental evidence about strategic uncertainty, focal points, and asymmetric coordination signals~\cite{vanhuyck1990,crawford2008,agranov2012}. Notebook outputs are conditional simulations; all human treatment effects remain hypotheses.

The scarce resource is differentiated transaction priority. This is not a full auction: choices are binary, and continuous bids, endogenous payments, and seller revenue are absent. Algorand fee and confirmation mappings, empirical payoffs, type probabilities, human sample size, and statistical tests must be calibrated or preregistered before a laboratory implementation.

\textbf{Computational scientist---GitHub:} \GitHubLink\\
\textbf{Runnable notebook:} \NotebookLink\\
\textbf{Behavioral scientist---Hugging Face:} \HuggingFaceLink\\
\textbf{Interdisciplinary scholar---A0 poster:} \PosterLink

\clearpage\balance
\section{Answer to Q3: behavior and interdisciplinary integration}
The rational benchmark assumes correct beliefs and payoff maximization. Alternatives include level-$k$ reasoning, focal attraction, ambiguity aversion, and reinforcement~\cite{camerer2003,zhu2025}. Treatment choices and elicited beliefs will distinguish them: convergence to type-dependent best responses supports the benchmark; persistent same-fee clustering after beliefs stabilize motivates revision. The current Hugging Face artifact is an exploratory behavioral interface with simulated co-players; no participant dataset currently exists.

\section{Advanced development, real-world impact, and roadmap}
The project connects Nash and Harsanyi foundations to a reproducible blockchain fee-choice experiment. Users, wallet designers, and protocol researchers could benefit from validated evidence about fee-menu or recommendation design. Risks are unrepresentative payoffs and treating a two-user laboratory game as a live-network model. Value therefore requires documented Algorand mapping, open code, robustness, and type-specific welfare checks. Academia requires identification and replication; industry adds operational tests; public institutions require accessibility and accountability. The next test is whether incomplete information causally raises same-fee congestion.

\begin{table}[h]\caption{From foundations to a collaborative interdisciplinary contribution.}
\begin{tabularx}{\columnwidth}{@{}p{1.6cm}Y@{}}\toprule
\textbf{Horizon}&\textbf{Milestone and evidence}\\\midrule
Foundation&Static $H/L$ game and equilibrium concepts.\\
PS2&Model, teaser, treatments, open code, and conditional simulated outputs; no empirical result.\\
Symposium&Planned poster, peer challenge, and response.\\
Next study&Calibration, preregistration, and human experiment.\\
2056&Accountable coordination tools for transaction markets.\\\bottomrule
\end{tabularx}\end{table}

\noindent\textbf{Expected 2056 Nobel Prize and Turing Award contribution statement.}
By 2056, we aspire to develop theory and auditable tools for coordination under private transaction urgency. Progress would require replicated experiments, field validation, open implementation, and independently verified welfare gains.

\subsection*{Open Science Statement}
Lin Zhang created the GitHub, notebook/Colab, and Hugging Face artifacts. The repository records URLs, license, parameters, seeds, commands, simulated outputs, and limitations. Zijie Wang will prepare the editable A0 poster.

\subsection*{Statement of Contribution to the UN's SDGs}
No SDG impact is demonstrated. A future SDG 9 claim requires evidence that the mechanism improves reliable access without disadvantaging patient or lower-value users.

\subsection*{Submission milestones}
Review-ready: September 27, 2026, 11:00 P.M. Beijing time. Symposium: September 28, 2026, 2:45--5:15 P.M., IB 2050. Final revision: September 30, 2026, 11:00 P.M. Beijing time.
